\section{Oblivious Multi-way Join with Filters}
\label{sec:bandjoin}

The residual query produced by \Cref{sec:decomposition} joins $k$ hop-relations
on their shared node ids and evaluates predicate filters on node
and edge properties. A natural extension of \wei to support filtered joins would treat
filtering as a post-processing step: materialize the full join output of
size $\mathit{OUT}_{\mathit{join}}$, then apply predicates to produce the
final result of size $\mathit{OUT}_{\mathit{filtered}}$. Schemes that
explicitly support filtering, such as
Obliviator~\cite{mavrogiannakis2025obliviator}, Opaque~\cite{zheng2017opaque}, and
ObliDB~\cite{eskandarian2019oblidb}, follow this
approach, revealing both sizes to the adversary and allocating memory
proportional to the unfiltered join output, \textit{which for selective
predicates can be orders of magnitude larger than the final result}.


\system avoids this by integrating filtering directly into the
multiplicity computation of \wei. Recall from
\Cref{sec:background:joins} that \wei associates every tuple with a final
multiplicity $\alpha_{\mathit{final}}$ --- the number of times it appears
in the output --- and that all downstream multiplicities are products. We
observe that a
tuple with $\alpha_{\mathit{final}} = 0$ also flows through every phase
identically to any other tuple: same memory accesses, same loop bounds,
same allocations. It simply expands to zero output rows in the expansion
phase, contributing nothing to the result.

\noindent\textbf{Filter as multiplicity.}
Each \onehop operator produces a hop relation of publicly known size $|E_e|$ in
which every row carries a boolean flag indicating whether that row passes
the query predicates (\Cref{sec:onehop}). When this hop relation enters
\wei as a leaf, we initialize its local multiplicity as:
\[
\alpha_{\mathit{local}}(t) =
\begin{cases}
1 & \text{if } t\text{'s filter flag is set,} \\
0 & \text{otherwise.}
\end{cases}
\]
Because all multiplicities propagate as products up and down the join
tree, a single $0$ at any tuple \textit{zeroes out its entire contribution to}
$\alpha_{\mathit{final}}$. A tuple that fails a filter never participates
in any output row, exactly as if it had been removed --- but it is never
actually removed, and no branch or allocation ever reflects its absence.
The adversary observes only $\mathit{OUT}_{\mathit{filtered}}$, the size
of the final filtered result, and never learns the size the join would
have produced without filtering.

\noindent\textbf{Correctness.} The modified initialization is the only
change to \wei's pipeline. All other phases --- bottom-up, top-down,
expansion, and alignment --- proceed identically. A tuple that fails a
filter has $\alpha_{\mathit{final}} = 0$ and expands to zero rows;
a tuple that passes has its multiplicity computed exactly as in the
unfiltered join. The output is therefore identical to first joining and
then filtering.

\noindent\textbf{Obliviousness.} Since $\alpha_{\mathit{local}}$ is
initialized from a flag stored inside the tuple rather than from any
branching on the tuple's value, the memory access pattern of every
phase remains a function of the input sizes and $\mathit{OUT}_{\mathit{filtered}}$
alone. No intermediate count
is allocated, computed, or revealed.

\noindent\textbf{Applicability.} The same technique applies to any
multiplicity-based oblivious join scheme. Obliviator~\cite{mavrogiannakis2025obliviator},
for instance, explicitly reveals both the join output size and the
filtered output size as separate public quantities. Replacing its
post-join filter with the multiplicity initialization above collapses
these two leakages into one, revealing only
$\mathit{OUT}_{\mathit{filtered}}$. The modification requires no
structural change to the underlying join algorithm.

% \section{The Oblivious Multi-Way Band Join, Step by Step}
% \label{sec:bandjoin}


% The residual query produced by the decomposition of \Cref{sec:decomposition} is
% evaluated by the oblivious multi-way join engine \wei
% (\Cref{sec:background:joins}). This section walks through that engine
% operationally, at the level of detail needed to follow the remainder of the
% paper. We first restate the one invariant the whole construction is built to
% preserve, then introduce the \emph{multiplicity} abstraction that drives it, and
% finally trace the four phases end to end on a small example. Throughout, we
% generalize from equality joins to \emph{band} joins, the form \system actually
% needs.

% %------------------------------------------------------------------------------
% \subsection{What ``oblivious'' must mean here}
% \label{sec:bandjoin:invariant}
% %------------------------------------------------------------------------------

% Under our threat model (\Cref{sec:background:threat}), the adversary observes the
% full sequence of memory accesses, the control flow, and the size of every
% allocation. The join is \emph{oblivious} when all of this observable behavior is
% a function of only three \emph{public} quantities:
% \begin{enumerate}[label=(\roman*),leftmargin=2.2em,itemsep=2pt]
%   \item the sizes of the input relations $|R_1|,\dots,|R_k|$;
%   \item the size of the output, $\mathit{OUT}$; and
%   \item the query itself --- its schema, predicates, and join tree.
% \end{enumerate}
% No branch, loop bound, array index, or buffer size may depend on a \emph{value}
% inside a tuple, nor on any intermediate count (such as how many tuples match a
% given key). The challenge the algorithm solves is to compute a join --- an
% inherently data-dependent operation --- while letting only these three public
% quantities show through. Every phase below meets this bar by replacing
% data-dependent control flow with two fixed-shape primitives: the \emph{oblivious
% sort} (osort), whose comparison network is fixed by the input length alone, and
% the \emph{linear scan}, a single left-to-right pass that applies the same
% branchless update to every adjacent pair of tuples.

% %------------------------------------------------------------------------------
% \subsection{The central abstraction: multiplicities}
% \label{sec:bandjoin:multiplicity}
% %------------------------------------------------------------------------------

% The engine never materializes an intermediate join. Instead it computes, for
% every input tuple $t$, a single integer: the number of rows of the \emph{final}
% result in which $t$ participates. We call this its final multiplicity
% $\alpha_{\mathit{final}}(t)$. Once these numbers are known, producing the output
% is mechanical --- copy each tuple $\alpha_{\mathit{final}}(t)$ times and line the
% copies up --- so the entire difficulty is reduced to computing multiplicities
% obliviously.

% Two facts make this abstraction powerful. First, $\alpha_{\mathit{final}}$
% factors along the join tree: a tuple's count in the whole join is its count
% within its own subtree times the count contributed by the rest of the tree.
% This is what lets the engine compute it in two sweeps --- one up the tree, one
% down. Second, because each output row contains exactly one tuple from each
% relation, the multiplicities of any single relation sum to $\mathit{OUT}$;
% summing $\alpha_{\mathit{final}}$ over a relation is how the engine learns the
% (public) output size without ever inspecting a value.

% \paragraph{Band predicates.}
% A child tuple $c$ matches a parent tuple $p$ not on equality but on a
% \emph{band}: $c.\mathit{key} \in [\,p.\mathit{key}+\delta_1,\ p.\mathit{key}+\delta_2\,]$,
% where $\delta_1,\delta_2$ are constants of the predicate (equality is the special
% case $\delta_1=\delta_2=0$, and one-sided inequalities use $\pm\infty$ bounds).
% The engine handles a band with a \emph{boundary} technique: for each parent
% tuple it emits two marker tuples, a \textsc{start} at $p.\mathit{key}+\delta_1$
% and an \textsc{end} at $p.\mathit{key}+\delta_2$. After an osort on the key,
% every child that matches $p$ falls between $p$'s two markers, so the number of
% matches is the difference of two running counts read off at the markers --- a
% quantity a single linear scan computes for all parents at once.

% %------------------------------------------------------------------------------
% \subsection{The four phases}
% \label{sec:bandjoin:phases}
% %------------------------------------------------------------------------------

% Given the join tree, the engine runs four phases in order. The first two compute
% multiplicities; the last two turn them into output. We describe each by its
% \emph{goal}, its \emph{mechanism}, and what it \emph{leaves behind}.

% \paragraph{Phase 1 --- Bottom-up: local multiplicities.}
% \emph{Goal.} Give every tuple $t$ its local multiplicity
% $\alpha_{\mathit{local}}(t)$: the number of partial results it participates in
% within the subtree rooted at its relation.
% \emph{Mechanism.} Every leaf tuple starts at $\alpha_{\mathit{local}}=1$. The
% engine then visits parent--child pairs in post-order. For each pair it
% concatenates the two relations (the parent expanded into \textsc{start}/\textsc{end}
% boundary tuples), osorts on the join key so matching children are bracketed by
% their parent's markers, and runs a linear scan that accumulates, for each parent
% tuple, the sum of its matching children's multiplicities. Multiplying in this
% contribution updates the parent's $\alpha_{\mathit{local}}$.
% \emph{Leaves behind.} After the sweep reaches the root, each tuple carries the
% size of the join of its own subtree. At the root specifically,
% $\alpha_{\mathit{local}}=\alpha_{\mathit{final}}$, since the root's subtree is the
% whole tree.

% \paragraph{Phase 2 --- Top-down: final multiplicities.}
% \emph{Goal.} Turn every tuple's local multiplicity into its final multiplicity
% $\alpha_{\mathit{final}}$ --- its count in the \emph{entire} join, not just its
% subtree.
% \emph{Mechanism.} Starting from the root (whose count is already final), the
% engine visits parent--child pairs in pre-order. Using the same
% concatenate--osort--scan machinery, but with the predicate read in the opposite
% direction, it pushes each parent's final count down to its children, distributing
% it across the matching children in proportion to their local multiplicities.
% \emph{Leaves behind.} Every tuple in every relation now carries
% $\alpha_{\mathit{final}}$, together with an offset that records where its copies
% belong in the output (used for alignment in Phase 4).

% \paragraph{Phase 3 --- Distribute-and-expand: replicate each tuple.}
% \emph{Goal.} Physically realize the counts: each relation must become a buffer in
% which tuple $t$ appears exactly $\alpha_{\mathit{final}}(t)$ times, padded to its
% output length.
% \emph{Mechanism.} A linear scan over a relation takes the prefix sum of
% $\alpha_{\mathit{final}}$, which gives both each tuple's destination offset and
% the relation's output size $\mathit{OUT}$. Tuples with count zero contribute
% nothing and are discarded; the survivors are scattered to their offsets by a
% fixed sequence of power-of-two ``distribute'' passes, leaving padding gaps
% between groups, and a final linear scan copies each tuple rightward to fill its
% gap. The number and stride of these passes depend only on $\mathit{OUT}$.
% \emph{Leaves behind.} Every relation in the tree now has exactly $\mathit{OUT}$
% rows, with each tuple occupying as many consecutive slots as its multiplicity.

% \paragraph{Phase 4 --- Align-and-merge: assemble the result.}
% \emph{Goal.} Glue the expanded relations together so that row $i$ holds a single,
% consistent combination of one tuple from each relation --- i.e.\ the $i$-th join
% result.
% \emph{Mechanism.} Walking the tree from the root, the engine osorts each child by
% an \emph{alignment key} derived from the offset recorded in Phase 2, so that the
% child's copies land in the same row order as the partial result they extend. It
% then concatenates the child's columns onto the accumulating result with a single
% lockstep pass over the two equal-length buffers.
% \emph{Leaves behind.} The accumulated relation at the root is the materialized
% join, of size $\mathit{OUT}$.

% %------------------------------------------------------------------------------
% \subsection{A worked example}
% \label{sec:bandjoin:example}
% %------------------------------------------------------------------------------

% Consider a two-relation equality join with parent $P=\{p_1,p_2\}$ and child
% $C=\{c_1,c_2,c_3\}$, where $p_1$ matches $c_1,c_2$ and $p_2$ matches $c_3$.

% \begin{itemize}[leftmargin=1.4em,itemsep=2pt]
%   \item \textbf{Phase 1.} Leaves $c_1,c_2,c_3$ start at
%         $\alpha_{\mathit{local}}=1$. The bottom-up scan over the $P$--$C$ pair
%         finds two matching children under $p_1$ and one under $p_2$, giving
%         $\alpha_{\mathit{local}}(p_1)=2$, $\alpha_{\mathit{local}}(p_2)=1$. As
%         root, these are already final.
%   \item \textbf{Phase 2.} Pushing the parents' counts down,
%         $\alpha_{\mathit{final}}(c_1)=\alpha_{\mathit{final}}(c_2)=
%         \alpha_{\mathit{final}}(c_3)=1$. Note both relations' multiplicities sum
%         to $3=\mathit{OUT}$.
%   \item \textbf{Phase 3.} $P$ expands to $[p_1,p_1,p_2]$ and $C$ to
%         $[c_1,c_2,c_3]$, each of length $3$.
%   \item \textbf{Phase 4.} Aligning $C$ against $P$ and concatenating yields the
%         three result rows $(p_1,c_1),(p_1,c_2),(p_2,c_3)$.
% \end{itemize}

% %------------------------------------------------------------------------------
% \subsection{Why the pipeline is oblivious}
% \label{sec:bandjoin:obliviousness}
% %------------------------------------------------------------------------------

% Every step above is one of three shapes: an osort whose network is fixed by the
% buffer length, a linear scan that touches every adjacent pair identically, or a
% distribute pass whose strides are fixed by $\mathit{OUT}$. None branches on a
% tuple's value, and none allocates a buffer sized by an intermediate count: the
% only sizes that ever appear are the input lengths, the per-relation output length
% $\mathit{OUT}$, and constants of the query. A tuple that matches nothing is not
% skipped --- it is carried along as a full-width record and silently dropped only
% when its zero multiplicity makes it padding in Phase 3. The pipeline therefore
% leaks exactly the three public quantities of
% \Cref{sec:bandjoin:invariant} and nothing more.

% %------------------------------------------------------------------------------
% \subsection{Oblivious Multi-Way Band Join with Filter}
% \label{sec:bandjoin:filter}
% %------------------------------------------------------------------------------

% The multi-way band join is driven entirely by multiplicities, not by
% deleting or keeping rows:
% \begin{itemize}[leftmargin=1.4em,itemsep=2pt]
%   \item \textbf{Bottom-up} computes each tuple's $\alpha_{\mathit{local}}$ (how
%         many times it matches within its subtree). We simply set
%         $\alpha_{\mathit{local}}=0$ for entries that do not pass the node filter.
%   \item \textbf{Top-down} combines these into $\alpha_{\mathit{final}}$ (the
%         total result rows the tuple participates in).
%   \item \textbf{Distribute-and-expand} replicates each tuple
%         $\alpha_{\mathit{final}}$ times.
%   \item \textbf{Align-and-merge} stitches the replicated tuples into the result.
% \end{itemize}

% A tuple with $\alpha_{\mathit{local}}=0$ flows through every phase identically:
% same memory accesses, same loop bounds, same allocations --- but because all
% downstream multiplicities are products, a single $0$ zeroes out its
% $\alpha_{\mathit{final}}$, so it expands to $0$ output rows. That is the entire
% trick: filtering is just zeroing a multiplicity, and the obliviousness falls out
% for free because nothing is skipped or resized.
